\section{Method}
\label{sec:method}

% =========================================================================
% Core section. Target length for first draft: ~5-6 pages.
% Mirrors design.md \S21 (Concrete Specialization), but written for a
% research-paper audience rather than as a handoff document.
% =========================================================================

\subsection{Overview}
\label{sec:method-overview}

The method is approximate $Q$-learning over a verifier-anchored context
MDP, with MCTS used as a policy improvement operator that supplies
stronger training targets during expert iteration and as a test-time
compute knob at inference. The central learnable object is a parametric
action-value function $\Qpsi(s, a)$ over structured context states; a
shared policy head $\ptheta(a \given s)$ supplies a categorical prior
for search and a sampling distribution for fast acting. The MDP is
defined by Sections~\ref{sec:method-state}--\ref{sec:method-actions}
(state representation and action set);
Section~\ref{sec:method-architecture} specifies the parametric heads
that approximate $\Qpsi$ and $\ptheta$;
Section~\ref{sec:method-reward} fixes the reward function;
Section~\ref{sec:method-training} states the approximate $Q$-learning
update with MCTS-supplied targets; Section~\ref{sec:method-mcts}
describes how MCTS produces the locally improved estimates $\Qmcts$ and
$\pimcts$ used as those targets; and Section~\ref{sec:method-inference}
states the two inference modes (fast $\Qpsi$ argmax, search-refined
MCTS) that the trained heads support.

\paragraph{What is borrowed and what is novel.}
The learning algorithm is approximate $Q$-learning
\citep{mnih2015human, sutton2018reinforcement}; the use of MCTS as a
policy improvement operator that produces training targets is expert
iteration \citep{anthony2017thinking, silver2017alphagozero,
silver2018alphazero}; and the bifurcation of the search tree into
decision and chance (afterstate) nodes, which lets the sampled chance
node estimate the inner Bellman expectation in a stochastic environment,
is borrowed from Stochastic MuZero
\citep{antonoglou2022stochastic}. The novel ingredients are
(i) the verifier-anchored context MDP itself
(Section~\ref{sec:method-state}); (ii) the shared-encoder, multi-head-$Q$
architecture, in which one frozen LLM serves as both encoder and
generator so a single forward yields the state embedding and the
verb-instantiator's continuation rather than running two models
(Section~\ref{sec:method-architecture}); and (iii) the
\textsc{Verify} action that makes checker calls a first-class
move within the MDP rather than a post-hoc filter
(Section~\ref{sec:method-actions}). The chance-node $Q$-estimate
$\Qmcts(s, a)$ is used as a direct Bellman target in place of
single-sample targets (Section~\ref{sec:method-mcts}); the afterstate
mechanism is borrowed, and what is specific to our setting is its
application to the non-enumerable LLM-induced transition, where the
sampled chance node is the only available estimator of that expectation.

\subsection{Structured Context State}
\label{sec:method-state}

The context MDP's state representation occupies a middle ground between
two failure modes. A state that is too unstructured (a raw token
sequence) gives the $Q$- and policy heads no inductive bias to exploit,
and makes any tree reuse impossible: two trajectories that arrive at
semantically equivalent situations never share state. A state that is
too compressed (a single learned embedding) loses interpretability and
cannot express the action-availability constraints that prevent the
agent from, say, invoking the verifier on a hypothesis that does not
exist. Our representation is a fixed structural shell whose slot
contents are LLM-generated text. The slots fix the kinds of entities the
agent can manipulate; the LLM fills in the specifics. The same
representation is consumed by the parametric heads
(Section~\ref{sec:method-architecture}), by the availability predicate
of Section~\ref{sec:method-actions}, and by MCTS when it samples
trajectories to refine local estimates (Section~\ref{sec:method-mcts}).

\begin{definition}[Math state]
\label{def:mathstate}
A \emph{math state} $s \in \Sset$ is a tuple
\[
  s = \bigl(\,P,\; F,\; G_{\textsc{open}},\; G_{\textsc{closed}},\;
            H,\; V,\; c,\; A,\; d\,\bigr),
\]
where $P \in \Sigma^{*}$ is the immutable problem statement, $F$ is a tuple
of \texttt{Fact} records, $G_{\textsc{open}}$ and $G_{\textsc{closed}}$ are
tuples of \texttt{Subgoal} records (open and closed), $H$ is a tuple of
\texttt{Hypothesis} records, $V$ is the verifier-call log, $c \in
\{\textsf{None}\} \cup \Sigma^{*}$ is the committed answer (if any), $A$ is
a tuple of identifiers naming explicitly abandoned hypotheses or subgoals,
and $d \in \mathbb{N}$ is the turn count.
\end{definition}

Concretely, each record type carries a stable identifier so that actions and
parent pointers can address it; the full \texttt{Fact}, \texttt{Subgoal},
\texttt{Hypothesis}, and \texttt{VerifierCall} dataclass definitions are
listed in Appendix~\ref{app:state-listings}.

\paragraph{The fact-versus-hypothesis distinction.}
The schema's most load-bearing element is the separation of \texttt{Fact}
from \texttt{Hypothesis}. A \texttt{Fact} is either a problem-given premise,
a derivation accepted as part of the reasoning trace, or a derivation that
has additionally been promoted by a successful \texttt{VerifierCall}. A
\texttt{Hypothesis} is a tentative claim that has not been (or could not be)
verified by the available checker. Without this distinction, the agent
cannot tell ``I have proven this'' from ``I think this is true,'' and the
\textsc{Verify} action (Section~\ref{sec:method-actions}) loses its
operational meaning. The distinction also exposes a structurally
interesting learnable behavior: spend a verifier call to promote a
hypothesis to a fact precisely when the chain of reasoning depends on it.

\paragraph{Recursive structure.}
The fields \texttt{parent\_hypothesis}, \texttt{children}, and
\texttt{follow\_up\_to} are forward-compatibility hooks for
\emph{structural decomposition}: a hypothesis may refine another
hypothesis, a subgoal may have nested sub-subgoals, and a verifier call may
follow up on an earlier one. The collections (\texttt{hypotheses},
\texttt{open\_subgoals}, \texttt{verifier\_log}) remain flat tuples; the
tree structure is implicit in the identifier references rather than
carried by nested data. Two consequences follow. First, the state encoder
(Section~\ref{sec:method-architecture}) does not require a graph neural
network or a tree-recursive cell; serializing the flat list with depth
markers (\texttt{[HYPOTHESIS id=hyp\_3 depth=2 parent=hyp\_1]}) lets the
frozen LLM's self-attention handle the tree natively. Second, the action
space stays categorical over a small fixed verb set
(Section~\ref{sec:method-actions}); the identifier of the action's target
is supplied by the verb-instantiator as part of its generated content.

We distinguish productive recursion (\emph{structural decomposition}: a
hypothesis splits into sub-hypotheses each addressing a sub-problem the
verifier can check) from unproductive recursion (\emph{meta-criticism}: a
critique of a critique, with no terminal verifier signal beneath). Phase 1
leaves the recursive fields at their defaults and so cannot produce
either; phase 2 activates them and relies on the verifier-anchored reward,
plus a per-hypothesis \texttt{tree\_depth\_limit} (default $3$ to $4$), to
suppress the unproductive case. Developing this distinction at depth,
and replacing the global depth limit with per-node decomposition
decisions, is deferred to a successor project on structured action
spaces.

\paragraph{Markov bet.}
The state is intended to satisfy the Markov property: all information
needed for future decisions lives in the slots, and the
parsed-from-LLM-text record is the agent's authoritative view. In
practice, the parser that maps LLM output into structured records is
lossy: an LLM sentence describing a derivation strategy may not survive
projection into \texttt{Fact.derivation}. The Markov claim is therefore a
design commitment and a debugging discipline rather than a theorem. When
a $Q$-value disagrees with expectation, the parser is the first suspect,
not the head.

\paragraph{Canonicalization.}
Two trajectories that reach semantically equivalent states (the same set
of facts and subgoals, discovered in different orders) typically yield
different schema instances under the natural trajectory ordering. MCTS
tree reuse benefits when such states hash identically, but
canonicalization requires a normalizer: topological ordering of facts by
derivation depth, SymPy normal forms on statements, sorted subgoal
identifiers. We defer canonicalization to phase 2; the first version
trades the reuse benefit for simpler debugging.

\subsection{Action Space}
\label{sec:method-actions}

The context MDP's action space confronts a similar tension to the state
representation. An action space that is too open (free-form text actions
generated directly by the LLM) gives the $Q$-head no fixed categorical
output set to score, breaks PUCT's reliance on a fixed action set, and
forces the policy head to score arbitrary candidates against arbitrary
candidates rather than learn a stable categorical distribution. An
action space that is too narrow (a rigid procedural workflow such as
\emph{classify-then-decompose-then-propose-then-critique-then-commit})
bakes a fixed deliberation order into the dynamics and prevents the
agent from learning to skip stages, revisit hypotheses, or invent
micro-loops. Our choice is a small fixed set of coarse verbs whose
targets and content are generated by the LLM as an action-conditioned
continuation. The verb gives the $Q$-head a small categorical output
set, the policy head a categorical distribution to learn, and MCTS
something categorical to branch on; the LLM-generated content lets the
action's effect adapt to the state.

\begin{definition}[Action: a verb applied to a target]
\label{def:actionset}
An action is a pair $a = (\textsc{verb}, \tau)$: a categorical \emph{verb}
drawn from a small fixed set, together with a \emph{target} identifier
$\tau$ naming the node of the state the verb operates on (absent for the
few verbs that introduce free-standing content). The six core verbs are
\textsc{Split}, \textsc{FormHypothesis}, \textsc{ComputeSubexpr},
\textsc{Verify}, \textsc{Commit}, and \textsc{Abandon}; an optional
reward-free \textsc{Critique} supports LLM self-assessment. The policy
factors accordingly into a categorical head over the verbs and a pointer
head over the live target identifiers; the verb-instantiator
(Section~\ref{sec:method-architecture}) supplies any free-form content as
a continuation of the serialized state. Each verb is a \emph{lens}: it
renders a scoped, grounded-local view of the target's neighborhood for the
frozen generator and writes the generated content back locally to that node
and its subtree.
\end{definition}

The verbs and their effects:

\begin{itemize}[leftmargin=*, itemsep=0.2em]
  \item \textsc{Split}: decompose a goal into children combined either
        conjunctively (subgoals, all of which must be solved) or, as an
        exhaustive case analysis, disjunctively (the cases together with an
        exhaustiveness obligation that is itself reified as a goal). The
        connective is recorded on the produced edge and consumed by
        \textsc{Commit}'s aggregation.
  \item \textsc{FormHypothesis}: introduce a tentative claim, a candidate
        solution addressing a goal.
  \item \textsc{ComputeSubexpr}: perform a symbolic or numeric computation
        step; the result enters as a \texttt{Fact} with
        \texttt{source = "derived"}.
  \item \textsc{Verify}: call the symbolic checker on a chosen hypothesis or
        derived fact; the call is appended to \texttt{verifier\_log} and, on
        success, the target is promoted to a verified fact. \textsc{Verify}
        is the only reward-bearing verb (Section~\ref{sec:method-reward}).
  \item \textsc{Commit}: close a goal by promoting a discharging witness (a
        verified fact, a solved decomposition, or a verified hypothesis) to
        its solution. \textsc{Commit} is recursive: committing the root goal
        is terminal and yields the verifier's verdict, while committing a
        subgoal is reward-free bookkeeping that records its solution for the
        parent.
  \item \textsc{Abandon}: close a goal or hypothesis negatively. This is a
        forward prune that marks the node abandoned and retains its subtree
        as a checked dead end, never a rollback to a prior state.
  \item \textsc{Critique} (optional, reward-free): an LLM self-assessment of
        a hypothesis against its goal. It adds information the policy may use
        to decide whether to commit, abandon, or refine, but never adds
        reward, so the symbolic verifier remains the sole reward source.
\end{itemize}

\textsc{Split} unifies conjunctive decomposition and case analysis under one
connective-tagged verb; compression of a closed subtree is handled by how the
lens renders it rather than by a verb. The \emph{Phase-1} implementation
realizes a conservative subset: a verb-only categorical action whose target
and content are sampled by the frozen generator from a global, unscoped view,
with verified-fact promotion and the terminal commit already in place. The
\emph{designed target} is the full $(\textsc{verb}, \tau)$ factorization with a
pointer head over identifiers and per-verb scoped lenses; the two are
distinguished wherever the difference matters below.

\paragraph{State-conditioned availability.}
Not every verb is meaningful in every state. We define an availability
predicate $\mathrm{avail}: \Sset \times \Aset \to \{0, 1\}$ that masks
invalid actions before PUCT:

\begin{center}
\small
\begin{tabular}{@{}ll@{}}
\toprule
Verb & Available when \\
\midrule
\textsc{Split}           & $|G_{\textsc{open}}| > 0$ and $d < d_{\max}$ \\
\textsc{FormHypothesis}  & $|G_{\textsc{open}}| > 0$ and $d < d_{\max}$ \\
\textsc{ComputeSubexpr}  & always \\
\textsc{Verify}          & $|H| > 0$ or $\exists f \in F:
                           f.\texttt{source} = \texttt{"derived"}$ \\
\textsc{Commit}          & a discharging witness exists for some open goal \\
\textsc{Abandon}         & $|H| > 0$ or $|G_{\textsc{open}}| > 0$ \\
\bottomrule
\end{tabular}
\end{center}

The state schema (Section~\ref{sec:method-state}) makes these predicates
trivial to evaluate. At each decision node, the policy logits for
unavailable verbs are masked to $-\infty$ before the softmax, and the
$Q$-head's unavailable entries are excluded from the PUCT argmax; under the
designed $(\textsc{verb}, \tau)$ factorization the same mask additionally
restricts the pointer head to type-valid targets.

\paragraph{The \textsc{Verify} design choice.}
\textsc{Verify} is the action that converts ``LLM formats output,
symbolic solver checks'' from a post-hoc parser pass into a first-class
deliberation move. The agent decides when to spend a verifier call;
intermediate results land in \texttt{verifier\_log} and become part of
the state. Three benefits follow. First, the verifier becomes a
learnable resource the policy can budget: a state in which the agent
has done substantial derivation work but has not invoked the verifier
carries implicit risk, and the policy can learn to call the verifier
when the expected information gain justifies the cost. Second, the
verifier log provides a structurally informative feature for the policy
head: trajectories with many successful intermediate verifier calls
correlate with eventual correctness. Third, the verifier can be invoked
on intermediate hypotheses rather than only the committed answer,
turning the binary terminal signal into a series of verifier-anchored
checkpoints, each of which can carry a small intermediate reward
(Section~\ref{sec:method-reward}) without introducing the Goodhart
concerns that attend learned-aesthetic signals.

\paragraph{The root \textsc{Commit} format constraint.}
The terminal root \textsc{Commit} requires strict-format output: the produced text
must parse for the terminal verifier (a boxed numeric expression for
GSM8K, a SymPy-evaluable expression for MATH, a Lean tactic for formal
proof tasks). Two implementations are possible:

\begin{enumerate}[leftmargin=*, itemsep=0.2em]
  \item \emph{Constrained decoding}: regex- or grammar-guided sampling
        \citep{willard2023outlines}. The format is guaranteed by
        construction; no retries.
  \item \emph{Retry with correction}: sample, parse, and on failure
        resample with a corrective prompt that includes the parse
        error.
\end{enumerate}

Path 1 is preferred when the chosen LLM and decoding framework support
it; path 2 is a portable fallback. The choice affects the
verb-instantiator (Section~\ref{sec:method-architecture}) but not the
$Q$-head or the action set: the root \textsc{Commit} is a single verb
regardless of how its content is produced.

\paragraph{Alternatives and migration path.}
We adopt the small fixed verb set as the conservative phase 1
starting point: the verb-instantiator supplies enough action-content
variation that MCTS still explores meaningfully, and the small fixed
action set keeps the multi-head $Q$ architecture
(Section~\ref{sec:method-architecture}) simple. Two alternative
action-space designs that recur in the LLM-MCTS literature (verb-times-slot
pairs and progressive widening with LLM-proposed candidates), and the
bounded migration path to the latter in phase 2, are discussed in
Appendix~\ref{app:action-alternatives}.

\subsection{Architecture: Shared LLM and Trainable Heads}
\label{sec:method-architecture}

The architecture of the approximate $Q$-learning system confronts a
third design tension. An external encoder (such as a sentence-transformer
trained on general text) introduces \emph{representation drift}: the
encoder's notion of state similarity need not match what the generator
considers the same context, so the $Q$-head can see two distinct states
where the generator would produce identical continuations. Full
fine-tuning of the LLM is expensive in compute and memory and conflates
two sources of improvement that we want to keep separate: the learning
loop and the raw representational capacity of the language model. The
middle path we adopt (Figure~\ref{fig:architecture}) uses one frozen LLM
as both encoder and generator, with a small task-tuned adapter and three
lightweight heads sharing the LLM's hidden state. Because encoder and
generator are the same model, representation drift is bounded by
construction; because the LLM is frozen, any measured improvement during
training is attributable to the heads and the search procedure.

\begin{figure}[t]
\centering
% Architecture figure for sec:method-architecture.
% Frozen LLM as shared encoder/generator; adapter MLP; three trainable heads.
%
% Conventions:
%   - Trainable components: light blue fill, solid border.
%   - Frozen components:    white fill, thick solid border, "(frozen)".
%   - Solid arrows:         data flow.
%   - Dashed arrow:         KV-cache reuse (LLM embedding pass to the
%                           verb-instantiator continuation).

\usetikzlibrary{positioning, arrows.meta, shapes.geometric, fit, calc,
                backgrounds}

\begin{tikzpicture}[
  font=\footnotesize,
  >={Stealth[length=2mm, width=1.4mm]},
  every node/.style={align=center},
  frozen/.style={
    draw, thick, rectangle, rounded corners=2pt,
    fill=white,
    inner sep=3pt
  },
  trainable/.style={
    draw, semithick, rectangle, rounded corners=2pt,
    fill=blue!8,
    inner sep=3pt
  },
  data/.style={
    draw, semithick, rectangle, rounded corners=1pt,
    fill=gray!8,
    inner sep=3pt
  },
  head/.style={
    draw, semithick, rectangle, rounded corners=2pt,
    fill=blue!8,
    minimum width=40mm,
    minimum height=10mm,
    inner sep=3pt
  },
  flow/.style={->, semithick},
  cache/.style={->, dashed, semithick, gray!45!black},
]

% =====================================================================
% Left column: input and serialized tokens (stacked vertically)
% =====================================================================
\node[data, minimum width=32mm, minimum height=16mm] (state) {%
  \textbf{MathState} $s$\\
  $(P, F, G_{\textsc{open}}, G_{\textsc{closed}},$\\
  $H, V, c, A, d)$
};

\node[data, below=12mm of state, minimum width=32mm, minimum height=16mm]
  (tokens) {%
  \textbf{Serialized tokens}\\
  \texttt{[PROBLEM]} \ldots\\
  \texttt{[FACT \ldots]} \ldots\\
  \texttt{[HYPOTHESIS]} \ldots\\
  \texttt{[STATE\_END]}
};

% =====================================================================
% Center column: LLM, then pool, then adapter (vertical stack)
% =====================================================================
\node[frozen, right=10mm of state,
      minimum width=32mm, minimum height=18mm] (llm) {%
  \textbf{Frozen LLM} \emph{(frozen)}\\
  causal self-attention\\
  hidden states $h_{1{:}T}$\\
  hidden dim $d$
};

\node[data, below=6mm of llm,
      minimum width=32mm, minimum height=10mm] (pool) {%
  Pool at \texttt{[STATE\_END]}\\
  output $h_T \in \R^{d}$
};

\node[trainable, below=6mm of pool,
      minimum width=32mm, minimum height=12mm] (adapter) {%
  \textbf{Adapter MLP}\\
  $\phi: \R^{d} \to \R^{m}$,\quad $m = 256$
};

% =====================================================================
% Right column: three heads, stacked vertically
% =====================================================================
\node[head, right=14mm of llm] (qhead) {%
  \textbf{$Q$-head} \quad $K = 7$ scalars\\
  \textsc{Decompose}, $\ldots$, \textsc{CommitAnswer}
};

\node[head, right=14mm of pool] (polhead) {%
  \textbf{Policy head} \quad $K$ logits\\
  masked by avail.\ $\to \softmax$
};

\node[head, right=14mm of adapter] (vihead) {%
  \textbf{Verb-instantiator}\\
  continuation from \texttt{[STATE\_END]}
};

% =====================================================================
% Arrows
% =====================================================================
% Left vertical flow
\draw[flow] (state) -- (tokens)
  node[midway, right=2pt]{$\mathrm{serialize}$};

% Tokens to LLM (sweeps right and up)
\draw[flow] (tokens.east) -- ++(6mm, 0) |- (llm.west)
  node[pos=0.75, above]{tokens};

% Central column vertical flow
\draw[flow] (llm.south) -- (pool.north);
\draw[flow] (pool.south) -- (adapter.north);

% Adapter fans right to all three heads via a vertical bus
\coordinate (fanBase) at ($(adapter.east) + (6mm, 0)$);
\draw[semithick] (adapter.east) -- (fanBase)
  node[midway, above]{$z \in \R^{256}$};
\draw[flow] (fanBase) |- (qhead.west);
\draw[flow] (fanBase) |- (polhead.west);
\draw[flow] (fanBase) |- (vihead.west);

% KV-cache reuse: dashed arrow from LLM directly to verb-instantiator.
% Routes above all three heads then drops down on the right-hand side,
% terminating at the right edge of the verb-instantiator box.
\coordinate (kvOut)  at ($(llm.north east) + (-3mm, 0)$);
\coordinate (kvTopL) at ($(kvOut) + (0, 5mm)$);
\coordinate (kvTopR) at ($(qhead.north east) + (2mm, 5mm)$);
\coordinate (kvDrop) at ($(vihead.east) + (2mm, 0)$);
\draw[cache]
  (kvOut)
  -- (kvTopL) -- (kvTopR)
  node[pos=0.5, above, font=\footnotesize\itshape, gray!40!black]
    {KV-cache reuse}
  |- (kvDrop)
  -- (vihead.east);

% =====================================================================
% Legend (bottom)
% =====================================================================
\node[draw, thin, rounded corners=2pt, fill=white,
      below=4mm of tokens.south, anchor=north,
      inner sep=4pt, font=\footnotesize] (legend) {%
  \begin{tabular}{@{}l@{\;}l@{\quad}l@{\;}l@{}}
    \tikz\draw[thick, fill=white] (0,0) rectangle (4mm,2.5mm); & frozen &
    \tikz\draw[semithick, fill=blue!8] (0,0) rectangle (4mm,2.5mm); & trainable\\[1pt]
    \tikz\draw[->, semithick, >={Stealth[length=1.5mm,width=1.2mm]}]
      (0,0) -- (5mm,0); & data flow &
    \tikz\draw[->, dashed, semithick, gray!45!black,
               >={Stealth[length=1.5mm,width=1.2mm]}]
      (0,0) -- (5mm,0); & KV-cache reuse\\
  \end{tabular}
};

\end{tikzpicture}

\caption{Architecture overview: frozen LLM as shared encoder and
  generator; adapter MLP; multi-head $Q$; policy head; verb-instantiator.
  Solid arrows show data flow; the dashed arrow indicates KV-cache reuse
  between the embedding forward and the verb-instantiator's continuation
  generation.}
\label{fig:architecture}
\end{figure}

\paragraph{Serialization.}
A math state $s$ is rendered into a token sequence by a deterministic
template using explicit structural markers (\texttt{[PROBLEM]},
\texttt{[FACT \ldots]}, \texttt{[HYPOTHESIS \ldots]},
\texttt{[VERIFIER\_CALL \ldots]}, and a closing \texttt{[STATE\_END]}
sentinel); the full template is given in
Appendix~\ref{app:state-listings}. Two properties of this template
matter. First, the structural markers
give the LLM's self-attention stable positional anchors, so per-slot
information is recoverable from the pooled hidden state even though no
slot-specific encoding is performed explicitly. Second, depth-indexed
identifiers (\texttt{depth=2 parent=hyp\_1}) render the implicit tree
structure from Section~\ref{sec:method-state} as a flat token sequence
that causal self-attention can resolve without a tree-aware encoder.
Abandoned entities are filtered out, and only the immutable problem
prefix is shared across steps within a trajectory.

\paragraph{Frozen LLM forward.}
Let $\mathrm{serialize}: \Sset \to \Sigma^{*}$ denote the template
function and $\mathrm{tokenize}: \Sigma^{*} \to \mathbb{N}^{T}$ the
LLM's tokenizer. The frozen LLM produces a hidden state at each token
position. We define the state embedding as the hidden state at the
position of the $\texttt{[STATE\_END]}$ sentinel:

\begin{equation}
\label{eq:embed}
\mathrm{embed}(s) \;=\; h_{T} \;\in\; \R^{d_h},
\qquad
h_{1{:}T} \;=\; \mathrm{LLM}\bigl(\mathrm{tokenize}(\mathrm{serialize}(s))\bigr),
\end{equation}

where $T$ is the index of $\texttt{[STATE\_END]}$ and $d_h$ is the LLM's
hidden dimension (typically $2048$ to $4096$ for the $1$-to-$7$ billion
parameter base models we consider). Causal self-attention guarantees
that $h_{T}$ has integrated the entire prefix, and decoder-only LLMs
are trained to make this position's hidden state informative for
next-token prediction. We exploit this directly: the same $h_{T}$ that
the LLM would use to predict the next token serves as our state
embedding for the heads. The immutable problem prefix is KV-cached
across steps within a trajectory, so the per-step incremental cost is
only the new tokens introduced by the latest action (a single new
hypothesis, fact, or verifier-call record, typically).

\paragraph{Adapter MLP.}
The LLM's hidden state $h \in \R^{d_h}$ is compressed by a small adapter:

\begin{equation}
\label{eq:adapter}
\phi: \R^{d_h} \to \R^{m},
\qquad
\phi(h) \;=\; W_{2}\,\sigma(W_{1} h + b_{1}) + b_{2},
\end{equation}

with $m = 256$ and $\sigma$ the GELU activation. The adapter does three
things. First, $m \ll d_h$ reduces head parameter counts by a factor of
$d_h / m \approx 10$, which matters during repeated MCTS rollouts. Second,
the adapter gives the heads a task-tuned representation without
unfreezing the LLM: it absorbs the projection that turns a
next-token-prediction hidden state into a useful $Q$-input. Third, the
adapter is the natural place to begin unfreezing in phase 2 (e.g.,
adding LoRA \citep{hu2022lora} on the LLM's last transformer layers),
since the adapter has already absorbed much of the task-specific
representational work and there is less for the LLM weights to learn.

\paragraph{Multi-head $Q$.}
The $Q$-head produces $K = 7$ scalar outputs in one forward pass:

\begin{equation}
\label{eq:qhead}
\Qpsi(s, \cdot) \;=\; W_{Q}\,\sigma\bigl(U_{Q}\,\phi(h)\bigr) + b_{Q} \;\in\; \R^{K},
\end{equation}

with $\Qpsi(s, a_{k})$ the $k$-th component. Backpropagation on a
transition $(s, a_{k}, r, s')$ updates only the parameters in row $k$
of $W_{Q}$ along with the shared body; the other six heads receive
zero gradient from that sample. Heads specialize without being
separate networks, and a single LLM forward yields all action values
needed at a decision node. At inference, fast mode
(Section~\ref{sec:method-inference}) returns
$\arg\max_{a \in \mathrm{avail}(s)} \Qpsi(s, a)$.

\paragraph{Policy head.}
The policy head shares the adapter with $Q$ and produces $K$ logits.
Letting $\mathrm{avail}(s) \subseteq \Aset$ denote the available
actions at $s$,

\begin{equation}
\label{eq:policy}
\ptheta(a \given s) \;=\;
\frac{\exp\bigl(\ell_{\theta}(s, a)\bigr) \cdot \mathbf{1}[a \in \mathrm{avail}(s)]}
     {\sum_{a' \in \mathrm{avail}(s)} \exp\bigl(\ell_{\theta}(s, a')\bigr)},
\end{equation}

with logits $\ell_{\theta}(s, \cdot) = W_{\pi}\,\sigma(U_{\pi}\,\phi(h))$.
The mask removes unavailable actions exactly, so the softmax normalizer
is computed only over feasible actions; assigning them probability zero
through a renormalization is mathematically equivalent to assigning
them logit $-\infty$ and is more numerically stable in practice.

\paragraph{Verb-instantiator.}
The verb-instantiator produces the free-form content and target
identifier of the chosen action. It is implemented as continuation
generation from the $\texttt{[STATE\_END]}$ sentinel: after the LLM
has consumed the serialized state, the next tokens it generates
(conditioned on a small verb-specific prompt prefix) become the
action's content. Because the verb-instantiator continues from the
same KV cache that produced $\mathrm{embed}(s)$, two outputs are
obtained from a single LLM forward pass: the hidden state for the
heads, and the next-token logits for content generation. The total
LLM compute per decision node is one forward, not two.

A small verb-specific prompt distinguishes between, say, ``produce a
decomposition of \texttt{hyp\_3}'' and ``invoke the verifier on
\texttt{hyp\_3} with verb \texttt{factor}.'' These prompts are fixed
templates parameterized by the chosen verb and the target identifier
selected by the policy.

\paragraph{Target $Q$ network.}
Approximate $Q$-learning is unstable when the Bellman target chases
its own moving estimate. Following the standard DQN remedy
\citep{mnih2015human}, we maintain a target network $Q_{\bar{\psi}}$
that is a frozen copy of $\Qpsi$ (and the adapter), synchronized
every $N_{\mathrm{sync}}$ optimizer steps. Bellman targets use $Q_{\bar{\psi}}$:

\begin{equation}
\label{eq:bellmantarget}
y \;=\; r + \gamma \max_{a' \in \mathrm{avail}(s')} Q_{\bar{\psi}}(s', a'),
\end{equation}

and the $Q$-loss is $\sum (\Qpsi(s, a) - y)^{2}$ over the replay
batch. The adapter is synchronized along with the $Q$-head because
both contribute to the Bellman target. The policy head is not part
of the target copy and trains directly toward the visit-count
distribution $\pimcts(\cdot \given s)$
(Section~\ref{sec:method-training}).

\paragraph{Phase 1 freezing rationale.}
In phase 1 the LLM is frozen entirely. Only the adapter, $Q$-head,
policy head, and the small verb-specific prompt prefixes train.
The total trainable parameter count is below five million across all
heads, two to three orders of magnitude smaller than the frozen LLM.
Three consequences worth flagging: optimizer state and memory
footprint are dominated by the heads rather than the LLM, so a
modest GPU suffices even with a $7$B-parameter base model; any
measured improvement during training is attributable to the heads
and the search procedure rather than capability gains in the LLM,
which is the comparison framing needed for
Section~\ref{sec:experimental-plan}; and when a $Q$-value disagrees
with expectation, the LLM is not in the suspect set (it is fixed),
narrowing debugging to the adapter, the heads, the serialization
template, or the parser. Phase 2 optionally unfreezes the last one
or two transformer layers via LoRA \citep{hu2022lora}; this is a
separate experiment whose contribution is measured against the
frozen baseline.

This architecture is realized end-to-end in the pilot: the shared-encoder
implementation runs the frozen base model in process on a $1.5$B substrate,
reads $\mathrm{embed}(s) = h_T$ from the same weights that generate verb
content, and confirms the one-forward-per-decision prefill in practice
(Section~\ref{sec:pilot-closedloop}).

\subsection{MCTS as a Policy Improvement Operator}
\label{sec:method-mcts}

The parametric $\Qpsi$ and $\ptheta$ from approximate $Q$-learning are
useful, but on a fresh state $s$ they encode only the average behavior
accumulated across the training distribution; they have not yet
allocated any compute specifically to $s$. MCTS produces locally
improved estimates by sampling the environment more intelligently than
the behavior policy would, using $\Qpsi(s, a)$ as a leaf evaluator and
$\ptheta(a \given s)$ as a PUCT prior to bias which trajectories get
expanded. The outputs of one search from $s$ are the locally refined
estimates $\Qmcts(s, a)$ and $\pimcts(a \given s)$ defined in
Equations~(\ref{eq:qmcts}) and~(\ref{eq:pimcts}) below. These serve two
purposes: during training they replace single-sample Bellman targets
with averages over many trajectories whose actions were chosen by
search rather than by the behavior policy (Section~\ref{sec:method-training});
at inference they refine the head estimate at a state on which extra
compute is being spent (Section~\ref{sec:method-inference}). The MCTS
variant we use sits between two extremes. The bifurcation of the tree
into decision and chance (afterstate) nodes is borrowed from Stochastic
MuZero \citep{antonoglou2022stochastic}, which introduced the afterstate
node as the device for planning with a learned model in a stochastic
environment; we adopt it rather than invent it. The bifurcation matters
because the chance-node mean value is exactly the empirical estimator of
the inner expectation
$\E_{s' \sim P(\cdot \given s, a)}[r + \gamma V^{\ast}(s')]$ in the
Bellman optimality equation, and collapsing decision and chance nodes
into a single node type would fold that statistic into the decision node
and discard the very quantity that converges to $Q^{\ast}(s, a)$. A fully
model-based MCTS that enumerates $P(s' \given s, a)$ exactly is also
unavailable to us: the LLM-induced transition is not enumerable, since
$s'$ is determined by the verb-instantiator's continuation tokens and
the parser that lifts those tokens into the schema of
Section~\ref{sec:method-state}. This non-enumerability is what makes the
borrowed mechanism specific to our setting: the sampled chance node is
the only available estimator of the Bellman inner expectation, since
neither enumeration nor a closed-form transition model is on offer. The
middle path is therefore chance-node bifurcation with sampled transitions
and per-edge reward backup, with explicit decision and chance nodes, the
LLM used only as a sampler of $s' \sim P(\cdot \given s, a)$, and returns
accumulated from the leaf inward so that each ancestor receives the return
computed from its own vantage point.

\paragraph{Node types.}
The search tree alternates between two node types:

\begin{itemize}[leftmargin=*, itemsep=0.2em]
  \item \emph{Decision nodes} hold a state $s$ and per-verb children
        (one chance node per action $a \in \Aset$ that the
        availability predicate admits). They record $N(s)$, the visit
        count.
  \item \emph{Chance nodes} hold a state-action pair $(s, a)$ and a
        list of child decision nodes (one per sampled next state). They
        record $N(s, a)$, the visit count; $W(s, a)$, the accumulated
        return; and the prior $\ptheta(a \given s)$.
\end{itemize}

The chance-node split is the structure that lets $W(s, a) / N(s, a)$
converge to the inner Bellman expectation as the number of simulations
grows.

\paragraph{The four-phase loop.}
Each simulation executes four phases.

\begin{enumerate}[leftmargin=*, itemsep=0.2em]
  \item \emph{Selection.} Starting at the root decision node, descend
        the tree by applying PUCT at each decision node and sampling a
        transition at each chance node. Continue until reaching either
        an unexpanded chance node, a leaf decision node, or a terminal
        state.
  \item \emph{Expansion.} If the descent stopped at a decision node
        that has not yet been expanded, instantiate one chance-node
        child per available action, query the heads
        (Section~\ref{sec:method-architecture}) once to obtain
        $\Qpsi(s, \cdot)$ and $\ptheta(\cdot \given s)$, and seed each
        new chance node with the prior.
  \item \emph{Simulation.} From the newly added decision node, obtain
        a leaf evaluation. Phase 1 uses the $V$-style estimate
        $\widehat{V}(s) = \max_{a \in \mathrm{avail}(s)} \Qpsi(s, a)$
        directly; no separate rollout policy is required, in keeping
        with the AlphaZero pattern \citep{silver2017alphagozero}. A
        rollout-based simulation is a phase-2 option if the head
        estimate proves too noisy.
  \item \emph{Backup.} Propagate the return from the leaf back through
        the path, accumulating per-edge rewards as it climbs. The
        precise update is described next.
\end{enumerate}

\paragraph{PUCT selection.}
At each decision node, the next action is chosen by the PUCT rule
\citep{silver2017alphagozero}, using the learned policy as prior and
masking unavailable actions:

\begin{equation}
\label{eq:puct}
a^{\ast}
\;=\;
\arg\max_{a \in \mathrm{avail}(s)}
\Bigl[\,
\widehat{Q}(s, a)
\;+\;
c_{\PUCT} \cdot \ptheta(a \given s) \cdot \frac{\sqrt{N(s)}}{1 + N(s, a)}
\,\Bigr],
\end{equation}

with $a \notin \mathrm{avail}(s)$ assigned PUCT score $-\infty$ before
the argmax. The exploitation term $\widehat{Q}(s, a)$ is the chance node's
empirical mean $W(s, a) / N(s, a)$ once the node has been visited and the
head estimate $\Qpsi(s, a)$ otherwise. The exploration term
$\ptheta(a \given s)$ comes from the
policy head and steers early-visit allocation toward actions the
policy distillation has identified as promising. The exploration
constant $c_{\PUCT}$ controls the relative weight; we take $c_{\PUCT}
\in [1.0, 1.5]$ as a starting range (Section~\ref{sec:method-training}).

\paragraph{Per-edge reward backup.}
For an MDP with intermediate rewards, the standard pattern of
``propagate one terminal scalar to every ancestor'' is incorrect even
when the leaf is terminal: each ancestor's chance node should accumulate
the return computed from its own position, including the discounted
contribution of every reward that follows on its branch. Let the path
visited during one simulation be a sequence of $(\text{chance}_t,
r_{t+1})$ pairs, where $\text{chance}_t$ is the chance node entered at
step $t$ and $r_{t+1}$ is the reward observed on the transition
$s_t \xrightarrow{a_t} s_{t+1}$. Let $\widehat{V}(s_L)$ be the leaf
evaluation. Returns are backed up by reverse accumulation:

\begin{equation}
\label{eq:backup}
G \;\gets\; \widehat{V}(s_L);
\qquad
\text{for } t = L-1, L-2, \ldots, 0:
\quad
G \;\gets\; r_{t+1} + \gamma\, G,
\qquad
\text{update } \text{chance}_t \text{ with } G.
\end{equation}

The chance-node update is
$W(s_t, a_t) \gets W(s_t, a_t) + G$ and
$N(s_t, a_t) \gets N(s_t, a_t) + 1$, with the parent decision node's
visit count incremented in step. This pattern is the obvious
discretization of the Bellman recursion: the inner expectation that
appears at $(s_t, a_t)$ is estimated by averaging returns $G$ that
themselves embed the discounted Bellman recursion further down the
branch. Collapsing the per-edge backup to a constant terminal scalar
biases every ancestor's $Q$-estimate by the difference between the
correct return and that terminal scalar.

\begin{remark}[Chance-node consistency]
\label{rem:mcts-consistency}
Under the UCT convergence conditions of \citet{kocsis2006bandit}
(sufficient exploration of every descendant, bounded rewards, and the
sample-size conditions required for the tail-bound argument), the
chance-node mean value $W(s, a) / N(s, a)$ estimates the inner expectation
$\E_{s' \sim P(\cdot \given s, a)}[\, r + \gamma \widehat{V}(s') \,]$ formed
from the leaf evaluator $\widehat{V}$; this coincides with the Bellman
expectation $\E_{s' \sim P(\cdot \given s, a)}[\, r + \gamma V^{\ast}(s') \,]$
only in the limit where the leaf evaluator is itself exact. Our
chance-node estimator inherits the same conditions; we do not re-derive
convergence here. Two practical caveats apply. First, MCTS rollout
estimators are biased at finite sample sizes: the leaf-value estimator
from random rollouts (or, in our setting, from $\widehat{V}(s_L)$ as
returned by the head $\Qpsi$) is consistent only in the limit, and at
the small simulation budgets typical of real deployment the chance-node
mean has nonzero bias that depends on the leaf evaluator's own bias and
on the selection rule's exploration regime. Second, the PUCT prior used
here is a deliberate departure from the original UCT bound and changes
the relevant constants; the adapted bound for PUCT-with-priors is what
the AlphaZero family relies on in practice but is not what the
\citet{kocsis2006bandit} theorem proves. We treat Section~\ref{sec:exp-ablations}
(the sanity-check ablations) as the operative test of the estimator's
practical adequacy on our verifiable-math tasks, rather than as a proof of the
asymptotic statement.
\end{remark}

\paragraph{Distillation targets.}
After $B$ simulations from a root state $s$, MCTS yields two improved
estimates that serve as supervised targets for the heads. The visit-count
policy is

\begin{equation}
\label{eq:pimcts}
\pimcts(a \given s)
\;=\;
\frac{N(s, a)^{1/\tau}}{\sum_{b \in \mathrm{avail}(s)} N(s, b)^{1/\tau}},
\end{equation}

with $\tau > 0$ a temperature parameter
(Section~\ref{sec:method-training}). The chance-node $Q$-estimate is

\begin{equation}
\label{eq:qmcts}
\Qmcts(s, a) \;=\; \frac{W(s, a)}{N(s, a)},
\qquad a \in \mathrm{avail}(s).
\end{equation}

Both targets come from the same MCTS rollout, so trajectory collection
is shared between policy and $Q$ training. The two targets are
complementary: $\Qmcts(s, a)$ is a scalar regression target for the
$Q$-head's approximate $Q$-learning update; $\pimcts(\cdot \given s)$
is a categorical distribution that the policy head matches via
cross-entropy (Section~\ref{sec:method-training}). Using
$\Qmcts(s, a)$ in place of the single-sample Bellman target
$r + \gamma \max_{a'} Q_{\bar\psi}(s', a')$ is what makes MCTS a policy
improvement operator in this method: the same approximate $Q$-learning
algorithm runs, but its targets average over a search-guided distribution
of trajectories rather than over the behavior policy. Distilling visit
counts back into the policy prior closes the expert-iteration loop in
the style of AlphaZero \citep{silver2017alphagozero, silver2018alphazero}.

\subsection{Reward Structure}
\label{sec:method-reward}

Reward design admits two familiar failure modes. A purely verifier-anchored
signal (a terminal binary call to a non-LLM checker, no intermediate
shaping) is objective but sparse, and on harder problems the agent
explores trajectories whose terminal reward is almost always zero,
yielding little learning signal per episode. A learned-aesthetic signal
(a reasoning-quality classifier, a deception detector, a stylistic
critic) is dense but introduces a Goodhart attack surface: the policy
will, over enough updates, learn to satisfy the classifier at the
classifier's resolution rather than to satisfy the underlying objective.
The middle path we adopt is to admit only verifier-anchored signals as
Bellman targets, and to reserve learned-aesthetic signals for the MCTS
prior, where they bias exploration but cannot rewrite the value
function.

\paragraph{Principle.}
A signal enters the reward channel (and thus a Bellman target) only if
it corresponds to an objective proposition checkable by a non-LLM
program: a unit test passes, SymPy reports equality between a
canonicalized answer and the ground truth, a Lean lemma type-checks, a
numerical equality holds within tolerance. Learned classifiers (a
reasoning-quality scorer trained on human ratings; a deception detector
trained on adversarial trajectories) do not qualify, regardless of how
well they correlate with human judgment in offline evaluation.

\paragraph{First-version reward (terminal-only).}
Phase 1 uses a sparse terminal reward and lets MCTS supply credit
assignment to earlier actions, in the AlphaZero pattern
\citep{silver2017alphagozero}:

\begin{equation}
\label{eq:rewardterminal}
r_T \;=\;
\begin{cases}
+1 & \text{if } \mathrm{verifier}(c, P) = \texttt{correct}, \\
\phantom{+}0 & \text{if depth limit reached without commit}, \\
-1 & \text{if } \mathrm{verifier}(c, P) = \texttt{incorrect},
\end{cases}
\qquad r_t = 0 \text{ for } t < T,
\end{equation}

with $c$ the committed answer (the contents of the
$\texttt{committed\_answer}$ slot when \textsc{Commit} fires) and
$P$ the problem statement. Per-edge backup (Equation~\ref{eq:backup})
distributes the terminal signal to ancestor chance nodes weighted by
$\gamma$; combined with PUCT-driven selection, this is sufficient to
produce nontrivial $Q$-estimates at intermediate states even with the
binary terminal signal.

\paragraph{Intermediate verifier-anchored rewards (fallback).}
If terminal sparsity becomes a blocker on harder benchmarks (small base
models that score near zero on AIME, say), we introduce small positive
rewards for verifier-anchored intermediate events. The simplest such
event is a successful \textsc{Verify} that promotes a
hypothesis to a verified fact: $r_t = +0.1$ on the transition that
sets $f.\texttt{source} = \texttt{"verified"}$. The reward remains
objective (the verifier's success bit is a deterministic non-LLM
output), additive, and bounded, and the cumulative bonus is capped by
the depth limit. Crucially, the intermediate signal does not introduce
the Goodhart attack surface attached to learned classifiers because
the underlying check is the actual objective, applied at a smaller
granularity.

\paragraph{The Goodhart argument.}
The asymmetry between verifier-anchored and learned-aesthetic signals
is structural, and it instantiates a result established formally in the
reward-hacking literature \citep{skalse2022reward, gao2023scaling}: an
optimizer run against a learned proxy converges on the proxy's defects,
and the gap between proxy and true objective widens under optimization
\citep{gao2023scaling}. In our terms, verifier-anchored signals close
the \emph{intrinsic} Goodhart gap by construction (the signal is a
deterministic program, not a learned proxy off whose training
distribution a policy can wander), while the \emph{extrinsic} gap (the
formally specified verifier failing to capture the full intended
objective) remains and is treated as a threat to validity
(Section~\ref{sec:exp-threats}) rather than assumed away. The full
intrinsic-versus-extrinsic taxonomy, with worked examples, is developed
in Section~\ref{sec:disc-verifier}, which this subsection defers to.

\paragraph{Channel decomposition deferred.}
The channel-decomposed reward sketched in the design notes
($R = w_{\text{correct}} R_{\text{correctness}} +
w_{\text{safe}} R_{\text{safety}} + \cdots$) is overkill for the
verifiable-math setting, where a single correctness channel suffices.
Channel decomposition becomes relevant when the domain expands to
open-ended deliberation, dialogue, or tool use, where ``the right
answer'' is no longer a single Boolean and reward must trade off
competing desiderata. That extension is out of scope for phase 1; the
present framework is structured so that adding channels later is a
local change to the reward function and a head-shape change on the
$Q$-network, not a redesign of the loop.

\subsection{Training Loop}
\label{sec:method-training}

The learning algorithm is approximate $Q$-learning
\citep{mnih2015human, sutton2018reinforcement} with MCTS-supplied
targets. The $Q$-loss is the standard approximate-$Q$-learning
mean-squared Bellman error; what is novel is the use of MCTS chance-node
estimates $\Qmcts(s, a)$ as the target in place of single-sample
Bellman targets $r + \gamma \max_{a'} Q_{\bar\psi}(s', a')$. The
policy distillation loss $\mathcal{L}_\pi$ is a complementary
expert-iteration step \citep{anthony2017thinking, silver2017alphagozero}
that turns the search-derived visit-count distribution $\pimcts$ into a
sharpened categorical prior, which in turn drives more informative
search in the next round. Each outer iteration therefore alternates
trajectory collection (via MCTS using the current heads as priors and
value estimates) and head updates (approximate $Q$-learning plus policy
distillation). Algorithm~\ref{alg:outer-loop} summarizes the loop; the
remainder of this subsection states the loss functions precisely,
describes the replay buffer, and lists the hyperparameters that matter
most in the verifiable-math setting.

\begin{algorithm}[t]
\caption{Outer training loop}
\label{alg:outer-loop}
\begin{algorithmic}[1]
\State Initialize $\Qpsi$, $\ptheta$, target $Q_{\bar\psi} \gets \Qpsi$
\For{iteration $= 1, 2, \ldots$}
  \State Collect trajectories via MCTS with $\Qpsi$, $\ptheta$ as priors
  \State Compute $\Qmcts(s, a)$ and $\pimcts(\cdot \given s)$ for each
         visited node
  \State Add samples to replay buffer
  \For{minibatch step $= 1, \ldots, M$}
    \State Sample $(s, a, r, s')$ batch
    \State $\mathcal{L}_Q \gets \sum (\Qpsi(s, a) - \Qmcts(s, a))^2$
    \State $\mathcal{L}_\pi \gets - \sum_s \sum_a \pimcts(a \given s)
           \log \ptheta(a \given s)$
    \State Step optimizer on $\mathcal{L}_Q + \lambda \mathcal{L}_\pi$
  \EndFor
  \State Every $N_{\mathrm{sync}}$ steps: $Q_{\bar\psi} \gets \Qpsi$
\EndFor
\end{algorithmic}
\end{algorithm}

\paragraph{Loss functions.}
The two heads train jointly against MCTS targets. Let $\mathcal{B}$
denote a minibatch sampled from the replay buffer; the $Q$-loss is
mean-squared error against the chance-node estimate
(Equation~\ref{eq:qmcts}),

\begin{equation}
\label{eq:lossq}
\mathcal{L}_Q(\psi)
\;=\;
\sum_{(s, a) \in \mathcal{B}}
\bigl(\, \Qpsi(s, a) - \Qmcts(s, a) \,\bigr)^{2},
\end{equation}

and the policy loss is the cross-entropy between the masked head
distribution (Equation~\ref{eq:policy}) and the visit-count
distribution (Equation~\ref{eq:pimcts}),

\begin{equation}
\label{eq:losspi}
\mathcal{L}_\pi(\theta)
\;=\;
-\sum_{s \in \mathcal{B}} \sum_{a \in \mathrm{avail}(s)}
\pimcts(a \given s)\, \log \ptheta(a \given s).
\end{equation}

The combined objective is

\begin{equation}
\label{eq:losstotal}
\mathcal{L}_{\text{total}}(\psi, \theta)
\;=\;
\mathcal{L}_Q(\psi) + \lambda\, \mathcal{L}_\pi(\theta),
\end{equation}

with $\lambda$ a scalar that balances the two heads (starting value
$\lambda = 1$; the policy loss is on a different scale than the $Q$
regression but in practice $\lambda \in [0.5, 2]$ has worked across
similar setups). The adapter receives gradient from both heads, so
$\psi$ and $\theta$ share representation through the adapter and
specialize only in their head-specific final layers.

\paragraph{Replay buffer.}
Collected transitions are stored in a replay buffer keyed by the
chance-node tuple $(s, a, r, s', \Qmcts(s, a),
\pimcts(\cdot \given s))$. Sampling priority is proportional to the
TD error $|\Qpsi(s, a) - \Qmcts(s, a)|$ at the time of last visit, in
the style of prioritized DQN. Buffer entries are evicted on a
first-in-first-out basis subject to a capacity cap; stale entries do
not break correctness because $\Qmcts$ values are stored with the
sample rather than recomputed from the current heads.

\paragraph{Phase 1 backprop scope.}
Gradients flow only through the adapter, the $Q$-head, the policy
head, and the verb-instantiator's verb-specific prompt prefixes; the
LLM is frozen. The freezing rationale and its compute-attribution
consequences are stated in Section~\ref{sec:method-architecture}.

\paragraph{Hyperparameters worth flagging.}
Three hyperparameters carry disproportionate weight in early
experiments. The discount factor $\gamma$ is set in $[0.99, 1.0]$;
because committed episodes are short (at most $d_{\max} \approx 20$
turns) and the only nonzero reward is terminal in phase 1, the
discounting effect is mild and $\gamma = 1$ is a defensible starting
choice that simplifies the per-edge backup. The PUCT exploration
constant $c_{\PUCT}$ starts in $[1.0, 1.5]$, with higher values when
the policy prior is still close to uniform (early training) and lower
values once the prior has sharpened. The visit-count temperature
$\tau$ in Equation~(\ref{eq:pimcts}) is annealed: high $\tau$ early
(near $\tau \approx 1.5$) smooths the policy target and encourages
exploration during distillation; later iterations anneal toward
$\tau = 1$ (proportional to raw visit counts) and then below $1$ for
sharper, more deterministic targets. The target-network sync interval
$N_{\mathrm{sync}}$ is on the order of one thousand optimizer steps; faster sync
destabilizes the Bellman target, slower sync wastes gradient signal.
We list these starting values as the operative reference; the final
configuration is tuned on the validation split of the smallest
benchmark before the headline runs begin.

This loop is now implemented: its closed-loop realization and honest
per-decision compute accounting are reported in
Section~\ref{sec:pilot-closedloop}, where the loop is the approximate
$Q$-learning of this subsection with MCTS supplying the targets.

\subsection{Inference Modes}
\label{sec:method-inference}

The same trained $\Qpsi$ and $\ptheta$ support two inference modes that
differ only in the search budget. A single nonnegative integer $B$
selects between them; $B = 0$ is the fast mode (the trained $Q$-policy
acted on directly) used for ordinary queries, and $B > 0$ is the search
mode (MCTS-refined inference, using the same $\Qpsi$ and $\ptheta$ as
leaf evaluator and PUCT prior) used when the caller is willing to spend
additional compute for higher accuracy. The two modes are not separate
algorithms; they are the same learned actor evaluated with different
amounts of search applied on top. This is the test-time compute knob:
the caller chooses where on the compute-versus-accuracy frontier to
operate, without retraining and without changing the model.

\paragraph{Fast mode ($B = 0$).}
Fast mode uses the $Q$-head directly. At each decision node, the next
action is the availability-masked argmax,

\begin{equation}
\label{eq:fastmode}
\pi_0(s) \;=\; \arg\max_{a \in \mathrm{avail}(s)} \Qpsi(s, a),
\end{equation}

so a single LLM forward per decision (one embedding pass plus one
verb-instantiator continuation) suffices. Latency and cost are
constant in the depth of deliberation, bounded only by $d_{\max}$ and
the per-step LLM forward.

\paragraph{Search mode ($B > 0$).}
Search mode runs PUCT-based MCTS for $B$ simulations from the current
state and samples (or argmaxes over) the resulting visit-count
distribution,

\begin{equation}
\label{eq:searchmode}
\pi_B(a \given s) \;=\; \pimcts(a \given s; B),
\end{equation}

with $\pimcts$ defined as in Equation~(\ref{eq:pimcts}). The $B$
parameter is the test-time compute knob: small $B$ refines the head
estimate locally; large $B$ performs extended deliberation. Reported
inference compute should include all LLM forwards executed during the
$B$ simulations, not only the root-level forwards, to make the
compute-versus-accuracy comparison honest
(Section~\ref{sec:experimental-plan}).

\paragraph{Why fast mode works at all.}
Fast mode is competitive precisely because the policy and $Q$-heads
were distilled from MCTS targets during training. If training has
converged, the $Q$-head approximates the chance-node Bellman estimate
$\Qmcts$ at the training-time simulation budget, and the policy head
approximates $\pimcts$. Fast mode therefore evaluates a policy that
is, in expectation, close to the search policy at the training-time
$B$, but at constant inference cost. The empirical question we put on
record (Section~\ref{sec:experimental-plan}) is how far this
approximation gets the model on each benchmark, and how much
additional headroom $B > 0$ opens at matched compute.

\paragraph{Persistent search trees (future extension).}
For anytime deliberation, the search tree built for one query can
prime the next. If a caller asks the system to ``think harder,'' the
existing chance-node statistics, sampled transitions, and leaf
evaluations remain valid as long as the root state is unchanged; an
additional $B'$ simulations refines the same tree rather than starting
over. Carrying state across turns of a multi-turn deliberation
(rerooting the tree on a child after a committed action; reusing
subtrees that survive the rerooting) is the natural generalization.
We list this as a future extension; phase 1 restarts the tree per
query, which keeps the implementation simpler and the comparison to
$B = 0$ cleaner.
